%! TEX root = main.tex

\begin{feedback}[\cref{ex:la-1}]
Parts (b) and (c) are correct.  In (b), your counterexample is sufficient: a subspace must be closed under scalar multiplication.  In (c), you correctly verify that the zero vector is present and that both required closure properties hold.

Part (a) has the right strategy and conclusion, but there is an error in the scalar-multiplication computation.  From $v_1=(x_1,y_1,z_1)$, one has
\[
  av_1=(ax_1,ay_1,az_1),
\]
not $(ax_1,ay_1,z_1)$.  Your following line uses $az_1$, so the intended argument is valid once the displayed vector is corrected.  Also, the final closure property is closure under \emph{vector addition}, not vector multiplication.
\end{feedback}

\begin{feedback}[\cref{ex:la-2}]
Your formulation of the span and the criterion $au+bv=w$ are correct.  The conclusion is not: $w$ \emph{does} lie in the span.  The arithmetic error occurs after substituting $a=4-3b$ into the second equation.  It should give
\[
  2(4-3b)+b=3,
\]
because the second equation is $2a+b=3$; you instead changed both the left expression and right-hand side.  Solving the first two equations correctly gives coefficients that also satisfy the third equation.  Recheck that final substitution and then state the corresponding linear combination of $u$ and $v$.
\end{feedback}

\begin{feedback}[\cref{ex:la-3}]
Correct.  Your row operations produce a row-reduced echelon form with pivots in columns $1$ and $3$.  You then correctly take columns $1$ and $3$ from the \emph{original} matrix, not the reduced matrix, as the basis.  The relations $c_2=2c_1$ and $c_4=2c_1+c_3$ justify that these two independent columns span the whole column space, so its dimension is $2$.
\end{feedback}

\begin{feedback}[\cref{ex:la-4}]
Correct: you exhibit a nontrivial linear combination of the three polynomials that equals the zero polynomial, so the list is linearly dependent.  The relevant condition is that the coefficients are \emph{not all zero}; they do not need to sum to zero.  Here the coefficients are $1,1,-1$, and the displayed polynomial identity is exactly the justification required by the definition.
\end{feedback}

\begin{feedback}[\cref{ex:la-5}]
The algebra is correct and it proves the result.  There is one logical issue in the initial statement: linear dependence means that there exist $a,b$ which are \emph{not both zero}, rather than that both $a$ and $b$ are nonzero.  Your calculation actually handles the correct, more general condition: it shows that any equation
\[
  a(u-v)+b(u+v)=0
\]
forces $a=b=0$.  State that directly as a proof of independence, or use ``not both zero'' in the contradiction setup.  The use of independence of $u,v$ to obtain $a+b=0$ and $b-a=0$ is exactly right.
\end{feedback}

\begin{feedback}[\cref{ex:la-6}]
Your first sentence correctly identifies redundancy: a column is redundant when it is a linear combination of the other feature columns.  This indeed implies $\rank A<n$ when $A$ has $n$ columns.  More precisely, it means there is a nonzero coefficient vector $x$ with $Ax=0$; saying that the matrix ``collapses'' vectors to the same line or plane is not the right general description.

The pivot-column claim needs one crucial correction.  The pivot columns of the \emph{row-reduced matrix} need not be a basis for the column space of $A$, because row operations generally change that column space.  Instead, row reduction tells you the pivot \emph{indices}; take the columns with those indices from the \emph{original matrix} $A$.  Those original columns form a nonredundant basis for the same column space.
\end{feedback}
